\documentclass[11pt,twocolumn,a4paper]{article}

\usepackage[utf8]{inputenc}
\usepackage[T1]{fontenc}
\usepackage{amsmath,amssymb,amsthm,mathtools}
\usepackage{bm}
\usepackage{physics}
\usepackage{graphicx}
\usepackage{hyperref}
\usepackage{cleveref}
\usepackage[margin=2.2cm]{geometry}
\usepackage{enumitem}
\usepackage{float}
\usepackage{booktabs}
\usepackage{natbib}
\usepackage{algorithm}
\usepackage{algorithmic}
\usepackage{listings}
\usepackage{xcolor}
\usepackage{caption}
\usepackage{microtype}
\usepackage{lmodern}
\usepackage{tikz}
\usetikzlibrary{arrows.meta,positioning,shapes.geometric,fit,calc}

% Theorem environments
\newtheorem{theorem}{Theorem}[section]
\newtheorem{lemma}[theorem]{Lemma}
\newtheorem{corollary}[theorem]{Corollary}
\newtheorem{proposition}[theorem]{Proposition}
\theoremstyle{definition}
\newtheorem{definition}[theorem]{Definition}
\newtheorem{axiom}{Axiom}
\newtheorem{construction}[theorem]{Construction}
\theoremstyle{remark}
\newtheorem{remark}[theorem]{Remark}
\newtheorem{example}[theorem]{Example}

% Custom commands
\newcommand{\kB}{k_{\mathrm{B}}}
\newcommand{\Sk}{S_{\mathrm{k}}}
\newcommand{\St}{S_{\mathrm{t}}}
\newcommand{\Se}{S_{\mathrm{e}}}
\newcommand{\Sspace}{\mathcal{S}}
\newcommand{\Scoord}{\mathbf{S}}
\newcommand{\R}{\mathbb{R}}
\newcommand{\Z}{\mathbb{Z}}
\newcommand{\N}{\mathbb{N}}
\newcommand{\C}{\mathbb{C}}
\newcommand{\Hilbert}{\mathcal{H}}
\newcommand{\Recv}{\mathcal{R}}
\newcommand{\Floor}{\mathfrak{S}_{\mathrm{floor}}}
\newcommand{\Osc}{\mathfrak{O}}
\newcommand{\Cat}{\mathfrak{C}}
\newcommand{\Part}{\mathfrak{P}}
\newcommand{\Reps}{\mathfrak{R}}
\newcommand{\Frame}{\mathfrak{F}}
\newcommand{\eval}{\mathrm{eval}}
\newcommand{\Sexp}{\xi}
\newcommand{\morphism}{\Phi}
\newcommand{\catalyst}{\mathcal{K}}
\newcommand{\Sig}{\Sigma}
\newcommand{\Leaf}{\mathcal{L}}
\newcommand{\Tree}{\mathcal{T}}

\hypersetup{
    colorlinks=true,
    linkcolor=blue,
    citecolor=blue,
    urlcolor=blue
}

\title{\textbf{An Expression Algebra for Biomolecules:\\
The Receiver $\Recv_{\mathrm{bio}}$ as a Type-Safe Morphism Chain over\\
Hierarchical Oscillator Leaves with Cross-Modal Conversion Functors}}

\author{
Kundai Farai Sachikonye\\[0.3em]
AIMe Registry for Artificial Intelligence\\
Experimental Bioinformatics, Maximus-von-Imhof-Forum 3\\
85354 Freising, Germany\\
\texttt{kundai.sachikonye@bitspark.com}
}

\date{\today}

\begin{document}
\maketitle

\begin{abstract}
We construct $\Recv_{\mathrm{bio}}$, a receiver in the sense of the unconstrained-substructures calculus, that evaluates S-expressions denoting biomolecular systems and returns categorical addresses, partition signatures, and oscillatory trajectories under a single typed morphism chain. The construction generalises the partition-calculus folding pipeline by promoting its three-view fusion to the recursive triple decomposition $(\Sexp_k, \Sexp_t, \Sexp_e)$ of the unconstrained-substructures calculus, and by extending the residue-oscillator leaf to a five-class leaf algebra encompassing residues, prosthetic cofactors, electronic modes, solvent oscillators, and substrate ligands. We establish: (i) a Floor Theorem $\Floor(\Recv_{\mathrm{bio}}) > 0$ with explicit estimate; (ii) Type Safety of the morphism chain $\eval_{\Recv_{\mathrm{bio}}} = \mathrm{access} \circ \mathrm{fuse} \circ \mathrm{catalyze}^* \circ \mathrm{observe}$ over heterogeneous leaves; (iii) closure of the conversion functors $\Phi^{R \to R'}$ between oscillatory, categorical, and partition representations under $S_3$ action; (iv) seven Instrument Identity Theorems establishing that each of the framework's measurement instruments---folding 2D-IR, hyperfine spectrometry, harmonic molecular resonance, ensemble strobes, multi-modal holograms, atoms-as-spectrometers, and electron trajectory tracing---is a receiver evaluation of a sub-expression of $\Sexp_{\mathrm{biomolecule}}$; (v) a principled $\tau$-assignment rule reducing the atoms-as-spectrometers ternary state to a deterministic function of the partition occupancy deviation under local equilibrium; (vi) a two-tier chirality resolution permitting empirical spin-crossover transitions while strictly preserving the topological selection rule $\Delta s_{\mathrm{orbital}} = 0$; (vii) a worked instantiation on the resting state of cytochrome CYP3A4 (PDB \texttt{1TQN}). The receiver provides the foundation for a fifteen-paper monograph completely characterising cytochrome P450 catalysis as a closed trajectory in S-entropy space.

\medskip
\noindent\textbf{Keywords:} unconstrained-substructures calculus, S-expression, receiver evaluation, partition coordinates, S-entropy coordinates, triple equivalence, conversion functor, morphism chain, oscillator leaf algebra, cofactor leaf, electronic leaf, solvent leaf, $\tau$-assignment, spin-crossover, cytochrome P450, CYP3A4
\end{abstract}

%==============================================================================
\part{Foundations}
%==============================================================================

%==============================================================================
\section{Introduction}
\label{sec:introduction}
%==============================================================================

\subsection{The Problem}

The framework comprising the categorical protein database \citep{sachikonye2025c}, the observation-computation equivalence \citep{sachikonye2025occ}, the partition-calculus folding pipeline \citep{sachikonye2025fold}, the spectroscopic derivation of elements \citep{sachikonye2025spec}, and the four physical instruments---categorical spectrometer \citep{sachikonye2025_hyperfine}, harmonic molecular resonator \citep{sachikonye2025_hmr}, ensemble strobes \citep{sachikonye2025_strobes}, superimposed multi-modal holograms \citep{sachikonye2025_holo}---is unified mathematically by the unconstrained-substructures calculus \citep{sachikonye2025unc}. Each component takes physical observations, generates oscillatory data, and returns either an address in S-entropy space $\Sspace = [0,1]^3$ or a partition signature in coordinates $(n, \ell, m, s)$.

In current practice each component carries its own evaluation procedure. The folding pipeline runs Kuramoto integration, Fourier-transforms the coupling time series, and threshold-classifies the resulting image \citep{sachikonye2025fold}. The categorical spectrometer enumerates partition cells and counts mode transitions \citep{sachikonye2025_hyperfine}. The harmonic resonator searches for rational frequency proximities and returns a graph cycle rank \citep{sachikonye2025_hmr}. These procedures share much underlying structure but their compatibility and composability has not been established.

This paper makes the compatibility explicit. We construct a single receiver $\Recv_{\mathrm{bio}}$ that evaluates an S-expression denoting any biomolecular system---polypeptide, prosthetic group, bound substrate, surrounding solvent, transferring electron---and returns the partition signature whose cells encode the biomolecular state. Each existing instrument becomes a sub-evaluation under $\Recv_{\mathrm{bio}}$ restricted to its appropriate sub-expression. The receiver provides the foundation against which the remaining fifteen monograph papers \citep{sachikonye2025_p450monograph} construct successively richer characterisations of cytochrome P450 catalysis.

\subsection{What This Paper Establishes}

We prove eight load-bearing results.

\begin{enumerate}[leftmargin=*]
\item \textbf{Floor Theorem for $\Recv_{\mathrm{bio}}$} (Theorem~\ref{thm:floor}): $\Floor(\Recv_{\mathrm{bio}}) > 0$, with explicit lower bound determined by the discretisation depth of the partition signature, the Q-factors of the hardware oscillators implementing the receiver, and the conversion functor floor errors.
\item \textbf{Type Safety} (Theorem~\ref{thm:type_safety}): the morphism chain $\eval_{\Recv_{\mathrm{bio}}} = \mathrm{access} \circ \mathrm{fuse} \circ \mathrm{catalyze}^* \circ \mathrm{observe}$ composes correctly across all five leaf classes.
\item \textbf{Conversion Closure} (Theorem~\ref{thm:cycle_closure}): the functors $\Phi^{R \to R'}$ between the three representations $\{\Osc, \Cat, \Part\}$ form a closed cycle under $S_3$ action, with composition $\Phi^{R' \to R''} \circ \Phi^{R \to R'} = \Phi^{R \to R''}$.
\item \textbf{S-Entropy Conservation} (Theorem~\ref{thm:s_conservation}): the morphism chain preserves S-entropy modulo bounded transformation kernels.
\item \textbf{Triple Decomposition Lifting} (Theorem~\ref{thm:triple_lift}): the partition-calculus folding pipeline's three-view fusion $(\Sig_{\mathrm{inst}}, \Sig_{\mathrm{avg}}, \Sig_{\mathrm{cat}})$ is the canonical recursive triple $(\Sexp_k, \Sexp_t, \Sexp_e)$ of the unconstrained-substructures calculus.
\item \textbf{Instrument Identity Schema} (Theorem~\ref{thm:instrument_schema}): each of the seven measurement instruments satisfies an identity $I_n(\text{output}) \propto |\eval_{\Recv_n}(\Sexp_n)|^2$ where $\Recv_n$ is the receiver restricted to the instrument's modality and $\Sexp_n$ is the appropriate sub-expression.
\item \textbf{$\tau$-Assignment Rule} (Theorem~\ref{thm:tau_rule}): the atoms-as-spectrometers ternary state $\tau \in \{0, 1, 2\}$ is determined by the sign of the partition occupancy deviation under the receiver's local-equilibrium prior, making every $\chi^2$ result reproducible from a structural file alone.
\item \textbf{Spin-Crossover Resolution} (Theorem~\ref{thm:spin_resolution}): the chirality coordinate splits into $s_{\mathrm{orbital}} \in \{-\tfrac{1}{2}, +\tfrac{1}{2}\}$ (strictly conserved by topological invariance) and $s_{\mathrm{state}}$ (effective magnetic projection, exchange-coupled, permitted to change). Spin-crossover transitions are receiver-internal physical reorganisations within a fixed categorical cell, not categorical transitions.
\end{enumerate}

\subsection{Paper Structure}

Part~I (Sections~\ref{sec:introduction}--\ref{sec:axioms}) restates the unconstrained-substructures calculus axioms in the form required for biomolecules. Part~II (Sections~\ref{sec:leaves}--\ref{sec:morphism_chain}) constructs the receiver: leaf algebra, the three representations, conversion functors, and the morphism chain. Part~III (Sections~\ref{sec:floor_proof}--\ref{sec:identity_theorems}) proves Floor, Type Safety, and the seven Instrument Identity Theorems. Part~IV (Sections~\ref{sec:tau_rule}--\ref{sec:miracle}) resolves the two foundational ambiguities ($\tau$-assignment and spin-crossover) and develops the Miracle Principle for catalytic intermediates. Part~V (Section~\ref{sec:cyp3a4}) instantiates the receiver on the resting state of CYP3A4 against the \texttt{1TQN} crystal structure. Part~VI (Sections~\ref{sec:limits}--\ref{sec:conclusion}) discusses limitations and future work.

%==============================================================================
\section{The Unconstrained-Substructures Calculus, Recapped}
\label{sec:axioms}
%==============================================================================

We restate the axioms and three load-bearing theorems of \citet{sachikonye2025unc} in the notation used here.

\subsection{S-Expressions and Receivers}

\begin{definition}[S-Expression]
\label{def:sexp}
An \emph{S-expression} $\Sexp$ is a finite syntactic object built from atoms (typed constants), operators (typed symbols), and parenthesised compositions. The set of all S-expressions in a typed signature $\Sigma$ is $\mathcal{E}(\Sigma)$. Each $\Sexp \in \mathcal{E}(\Sigma)$ is tagged with a representation label $R(\Sexp) \in \Reps$, where $\Reps = \{\Osc, \Cat, \Part\}$ denotes the three permitted representations: oscillatory ($\Osc$), categorical ($\Cat$), and partition ($\Part$).
\end{definition}

\begin{definition}[Receiver]
\label{def:receiver}
A \emph{receiver} $\Recv$ is a tuple $(\Hilbert_\Recv, \eval_\Recv, \{\Phi_\Recv^{R \to R'}\}_{R, R' \in \Reps}, \Floor(\Recv))$ where:
\begin{itemize}[nosep]
    \item $\Hilbert_\Recv$ is the receiver's representation Hilbert space,
    \item $\eval_\Recv: \mathcal{E}(\Sigma) \to \Hilbert_\Recv$ is the evaluation map,
    \item $\Phi_\Recv^{R \to R'}: \mathcal{E}(\Sigma)|_R \to \mathcal{E}(\Sigma)|_{R'}$ are conversion functors between representations,
    \item $\Floor(\Recv) > 0$ is the receiver's intrinsic floor.
\end{itemize}
\end{definition}

A receiver is \emph{representation-coherent} if $\eval_\Recv \circ \Phi_\Recv^{R \to R'} = \eval_\Recv$ for all $R, R' \in \Reps$ and all expressions, modulo $\Floor(\Recv)$. The framework's Triple Equivalence \citep{sachikonye2025_partition} guarantees that representation-coherent receivers exist.

\subsection{The Three Load-Bearing Theorems}

\begin{theorem}[Floor; \citealt{sachikonye2025unc}, Theorem~T1]
\label{thm:floor_unc}
For every bounded receiver $\Recv$, $\Floor(\Recv) > 0$. Perfect alignment is impossible; the smallest knowable error is receiver-intrinsic.
\end{theorem}

\begin{theorem}[Triple Equivalence; \citealt{sachikonye2025_partition}, Theorem~3.1]
\label{thm:triple_eq}
The three representation categories are equivalent:
\begin{equation}
\mathbf{Osc} \simeq \mathbf{Cat} \simeq \mathbf{Part}
\end{equation}
with mutually-inverse functor pairs $(F_{OC}, F_{CO})$, $(F_{CP}, F_{PC})$, $(F_{PO}, F_{OP})$. All three categories yield identical entropy $S = \kB M \ln n$ for $M$ active modes at principal depth $n$.
\end{theorem}

\begin{theorem}[Unconstrained Subtask; \citealt{sachikonye2025unc}, Theorem~T3]
\label{thm:unconstrained}
For an S-expression $\Sexp = \Sexp_1 \star \cdots \star \Sexp_n$ evaluated by a representation-coherent receiver, the global value $\eval_\Recv(\Sexp) = k^*$ imposes no constraint on the local values $\eval_\Recv(\Sexp_i)$. Sub-expressions may be locally maximally infeasible (S = 100, ``miracles'') and the global S still equals $\Floor(\Recv)$.
\end{theorem}

These three theorems---Floor, Triple Equivalence, and Unconstrained Subtask---are the only foundational results we will assume. All other claims in this paper are derived.

\subsection{Recursive Triple Decomposition}

\begin{definition}[Recursive Triple]
\label{def:recursive_triple}
Every S-expression $\Sexp$ admits a recursive decomposition
\begin{equation}
\Sexp = (\Sexp_k, \Sexp_t, \Sexp_e)
\end{equation}
with $\eval_\Recv((\Sexp_k, \Sexp_t, \Sexp_e)) = \eval_\Recv(\Sexp)$ modulo $\Floor(\Recv)$. The components correspond to the three S-entropy axes:
\begin{itemize}[nosep]
    \item $\Sexp_k$: knowledge axis ($\Sk$, categorical structure, type information),
    \item $\Sexp_t$: temporal axis ($\St$, oscillatory dynamics, frequency content),
    \item $\Sexp_e$: evolution axis ($\Se$, partition refinement, energy distribution).
\end{itemize}
Recursing this decomposition gives a tree where each node has three children labelled $\{k, t, e\}$. Coordinate paths $\bm{p} \in \{k, t, e\}^d$ index sub-expressions at recursion depth $d$.
\end{definition}

This recursive triple is the link between the calculus and the partition-calculus folding pipeline. The latter's three-view fusion $(\Sig_{\mathrm{inst}}, \Sig_{\mathrm{avg}}, \Sig_{\mathrm{cat}})$ is---we will prove (Theorem~\ref{thm:triple_lift})---the canonical instance of $(\Sexp_k, \Sexp_t, \Sexp_e)$.

\subsection{The Bounded Phase Space Axiom}

\begin{axiom}[Bounded Phase Space; \citealt{sachikonye2025spec}, Axiom~2.1]
\label{ax:bounded}
All persistent dynamical systems occupy bounded regions of phase space with finite Liouville measure, and these bounded regions admit hierarchical partitioning into distinguishable subregions.
\end{axiom}

For biomolecules this axiom is empirically motivated: a folded protein occupies a finite volume bounded by the polypeptide chain's length and the fold's compactness; a prosthetic cofactor is bounded by its coordination sphere; an electron in a $d$-orbital is bounded by the Coulomb potential of its nucleus.

\begin{theorem}[Oscillatory Necessity; \citealt{sachikonye2025spec}, Theorem~2.4]
\label{thm:osc_nec}
For self-consistent dynamical systems in bounded phase space, oscillatory dynamics is the unique valid mode. Static, monotonic, and chaotic alternatives violate fundamental consistency requirements (Poincar\'e recurrence, finite information, reproducibility).
\end{theorem}

This converts every biomolecular component into a bounded oscillator, which is the starting point of the leaf algebra (Section~\ref{sec:leaves}).

\subsection{The Partition Coordinate System}

\begin{theorem}[Partition Coordinates; \citealt{sachikonye2025spec}, Theorem~3.4]
\label{thm:partition_coords}
A hierarchically nested oscillatory system admits coordinates $(n, \ell, m, s)$ with:
\begin{align}
n &\in \N^+ && \text{(principal partition depth)} \\
\ell &\in \{0, 1, \ldots, n-1\} && \text{(angular complexity)} \\
m &\in \{-\ell, \ldots, +\ell\} && \text{(orientation)} \\
s &\in \{-\tfrac{1}{2}, +\tfrac{1}{2}\} && \text{(chirality)}
\end{align}
The capacity at depth $n$ is $C(n) = 2n^2$, and transitions satisfy the selection rules
\begin{equation}
\Delta\ell = \pm 1, \quad |\Delta m| \leq 1, \quad \Delta s = 0.
\label{eq:selection_rules}
\end{equation}
\end{theorem}

The selection rule $\Delta s = 0$ is topological in origin (from $\pi_1(SO(3)) = \Z_2$). It is the cause of the spin-crossover ambiguity that we resolve in Section~\ref{sec:spin}.

%==============================================================================
\part{Constructing the Receiver}
%==============================================================================

%==============================================================================
\section{The Oscillator Leaf Algebra}
\label{sec:leaves}
%==============================================================================

\subsection{Five Classes of Leaves}

\begin{definition}[Oscillator Leaf]
\label{def:leaf}
An \emph{oscillator leaf} is a tuple
\begin{equation}
\Leaf = (\omega, \Scoord, \phi, R, \kappa)
\end{equation}
where $\omega \in \R_{>0}$ is the natural frequency, $\Scoord = (\Sk, \St, \Se) \in \Sspace$ is the S-entropy coordinate, $\phi \in [0, 2\pi)$ is the instantaneous phase, $R \in \Reps$ is the representation tag, and $\kappa$ is the leaf class:
\begin{equation}
\kappa \in \{\mathrm{res}, \mathrm{cof}, \mathrm{ele}, \mathrm{sol}, \mathrm{sub}\}.
\end{equation}
\end{definition}

The five leaf classes are:

\textbf{Residue} ($\kappa = \mathrm{res}$): an amino acid in the polypeptide chain, with $\omega$ derived from the amide-I band ($\sim$1650~cm$^{-1}$) shifted by side-chain identity, and S-entropy from Kyte--Doolittle hydrophobicity, van der Waals volume, and electrostatic index \citep{kyte1982, sachikonye2025c}.

\textbf{Cofactor} ($\kappa = \mathrm{cof}$): a prosthetic group such as heme b, FAD, FMN, NADPH, or substrate dioxygen. Cofactor leaves admit a nested decomposition---the heme leaf contains a porphyrin sub-leaf (HMR cycle rank), an Fe sub-leaf (Categorical Spectrometer at $Z = 26$), and an axial ligand sub-leaf (S 3p partition coordinates for thiolate).

\textbf{Electronic} ($\kappa = \mathrm{ele}$): an electron occupying a partition cell, with $\omega$ at the corresponding atomic transition frequency. The electronic leaf is the bridge to the spectroscopic-derivation paper \citep{sachikonye2025spec}: each electron is identified with a partition coordinate $(n, \ell, m, s)$ in its host atom, and the categorical-physical commutation $[\hat{O}_{\mathrm{cat}}, \hat{O}_{\mathrm{phys}}] = 0$ permits its trajectory to be tracked under categorical observation \citep{sachikonye2025_azurin, sachikonye2025_sod1}.

\textbf{Solvent} ($\kappa = \mathrm{sol}$): an ordered water molecule or counter-ion in the protein's hydration shell. The active-site water network, particularly the I-helix water cluster of cytochrome P450, plays a catalytic role through proton delivery; its inclusion as an explicit leaf class rather than implicit medium is essential for treating proton-coupled electron transfer.

\textbf{Substrate} ($\kappa = \mathrm{sub}$): a small-molecule ligand bound in the active site. Substrate leaves are S-entropy-coordinate-tagged like residue leaves but with $\omega$ and $\Scoord$ derived from the substrate's vibrational and electronic structure rather than from the canonical amino acid table.

\subsection{Leaf-Type Frequency Assignments}

For each leaf class we specify the frequency derivation rule.

\begin{definition}[Residue Frequency]
\label{def:res_freq}
For residue leaf with amino-acid identity $a$,
\begin{equation}
\omega_{\mathrm{res}}(a) = \omega_{\mathrm{amide-I}}^0 + \Delta\omega(a),
\end{equation}
where $\omega_{\mathrm{amide-I}}^0 = 1650~\text{cm}^{-1}$ and $\Delta\omega(a)$ is the side-chain-induced shift, tabulated in \citep{barth2007} and ranging from $-30$ to $+25$~cm$^{-1}$.
\end{definition}

\begin{definition}[Cofactor Frequency]
\label{def:cof_freq}
For cofactor leaf with cofactor identity $c$, $\omega_{\mathrm{cof}}(c)$ is the dominant vibrational mode of $c$. Specifically:
\begin{itemize}[nosep]
    \item Heme b: porphyrin breathing $\nu_4 \approx 1370$~cm$^{-1}$ (oxidation marker), $\nu_2 \approx 1580$~cm$^{-1}$ (spin marker), Fe-N stretch $\approx 240$~cm$^{-1}$.
    \item FAD/FMN: isoalloxazine ring modes $\sim$1580, 1410, 1350~cm$^{-1}$ \citep{abe2002}.
    \item NADPH: nicotinamide ring modes $\sim$1690, 1630, 1340~cm$^{-1}$.
    \item O$_2$ bound to Fe: $\nu(\mathrm{O\text{-}O}) \approx 1140$~cm$^{-1}$ (peroxo) to 800~cm$^{-1}$ (oxo) depending on protonation/oxidation state.
\end{itemize}
\end{definition}

\begin{definition}[Electronic Frequency]
\label{def:ele_freq}
For electronic leaf at partition coordinate $(n, \ell, m, s)$,
\begin{equation}
\omega_{\mathrm{ele}}(n, \ell, m, s) = \frac{|E_n - E_{n-1}|}{\hbar} = \frac{(2n - 1) Z_{\mathrm{eff}}^2 R_\infty c}{\hbar n^2 (n-1)^2}.
\end{equation}
For Fe $3d$ electrons in heme: $\omega_{\mathrm{ele}} \sim 10^{14}$~Hz (visible/near-IR), corresponding to the d-d transitions in the Q-band region.
\end{definition}

\begin{definition}[Solvent Frequency]
\label{def:sol_freq}
For solvent leaf, $\omega_{\mathrm{sol}} \approx 3400$~cm$^{-1}$ (O-H stretch) for water, with a 100--200~cm$^{-1}$ shift depending on hydrogen-bonding environment in the active site.
\end{definition}

\begin{definition}[Substrate Frequency]
\label{def:sub_freq}
For substrate leaf with chemical identity $\sigma$, $\omega_{\mathrm{sub}}(\sigma)$ is the substrate's reactive-bond stretching frequency---typically C-H ($\sim$2900~cm$^{-1}$) for hydroxylation substrates, C=C ($\sim$1640~cm$^{-1}$) for epoxidation substrates.
\end{definition}

\subsection{The Receiver Tree}

\begin{definition}[Receiver Tree]
\label{def:receiver_tree}
The receiver tree $\Tree(\Sexp_{\mathrm{biomolecule}})$ is a finite ordered tree whose interior nodes are tagged with composition operators ($\star$, $\diamond$, $+$) and whose leaves are oscillator leaves. The tree's root is the whole biomolecule expression. Recursion at each interior node follows Definition~\ref{def:recursive_triple}, producing three children indexed by $\{k, t, e\}$.
\end{definition}

\begin{example}[CYP3A4 Resting State Receiver Tree]
\label{ex:cyp3a4_tree}
The S-expression for CYP3A4 in its resting state $(\mathrm{Fe}^{3+}\text{-H}_2\mathrm{O}, \text{no substrate})$ decomposes as
\begin{align*}
\Sexp_{\mathrm{CYP3A4,rest}} &= (\Sexp_k^{\mathrm{CYP3A4,rest}}, \Sexp_t^{\mathrm{CYP3A4,rest}}, \Sexp_e^{\mathrm{CYP3A4,rest}})
\end{align*}
where
\begin{align*}
\Sexp_k^{\mathrm{CYP3A4,rest}} &= \mathrm{polypeptide(CYP3A4)} \\
&\quad \diamond \mathrm{cofactor}(\mathrm{heme}_{\mathrm{b}}) \\
\Sexp_t^{\mathrm{CYP3A4,rest}} &= \mathrm{Kuramoto}(\mathrm{H\text{-}bond\ network}) \\
\Sexp_e^{\mathrm{CYP3A4,rest}} &= \mathrm{partition}(\mathrm{Fe}\ 3d^5,\ \mathrm{HS},\ S=5/2) \\
&\quad \diamond \mathrm{partition}(\mathrm{H_2O\ axial})
\end{align*}
At depth one the tree has three children: a polypeptide-and-cofactor sub-tree (knowledge axis), a Kuramoto-network sub-tree (temporal axis), and an electronic-partition sub-tree (evolution axis). Each child recurses to yield individual leaves. The full tree for CYP3A4 has $\sim 10^4$ leaves.
\end{example}

%==============================================================================
\section{The Three Representations and Conversion Functors}
\label{sec:representations}
%==============================================================================

\subsection{Representation Categories}

\begin{definition}[Oscillatory Representation]
\label{def:osc_rep}
$\Osc$-representation: an S-expression in $\Osc$ specifies the system as a finite collection of oscillator leaves with explicit $(\omega_i, \Scoord_i, \phi_i)$ values, coupled through a coupling tensor $K_{ij}$ derived from S-entropy distance:
\begin{equation}
K_{ij} = K_0 \cdot f(d_\Sspace(\Leaf_i, \Leaf_j)) \cdot g(\rho(\Leaf_i, \Leaf_j))
\label{eq:coupling_general}
\end{equation}
where $f$ is a Gaussian kernel (Equation~\eqref{eq:coupling_kernel} below), $\rho$ is a class-aware separation function generalising the residue sequence separation, and $g$ is the corresponding decay factor.
\end{definition}

\begin{definition}[Categorical Representation]
\label{def:cat_rep}
$\Cat$-representation: an S-expression in $\Cat$ specifies the system as an object in a category with morphisms preserving the partition structure. Each leaf is assigned a triple categorical label $(\lfloor \log_2 \omega \rfloor, \lfloor (\omega - 2^{\lfloor \log_2 \omega \rfloor}) / 2^{\lfloor \log_2 \omega \rfloor - 1} \rfloor, \arg \Phi(\Leaf, 0))$ encoding (octave, fractional bit, phase) \citep{sachikonye2025_partition, sachikonye2025_hyperfine}.
\end{definition}

\begin{definition}[Partition Representation]
\label{def:part_rep}
$\Part$-representation: an S-expression in $\Part$ specifies the system as a hierarchical partition of bounded phase space, with each leaf assigned coordinates $(n, \ell, m, s)$ obeying the capacity formula $C(n) = 2n^2$ and the selection rules $\Delta\ell = \pm 1$, $|\Delta m| \leq 1$, $\Delta s_{\mathrm{orbital}} = 0$.
\end{definition}

The coupling kernel $f$ in Equation~\eqref{eq:coupling_general} is, for all classes,
\begin{equation}
f(d) = \exp\!\left( - \frac{d^2}{2\sigma_\kappa^2} \right)
\label{eq:coupling_kernel}
\end{equation}
with class-dependent width $\sigma_\kappa$. For residue-residue $\sigma_{\mathrm{res,res}} = 0.30$, for residue-cofactor $\sigma_{\mathrm{res,cof}} = 0.20$, for cofactor-cofactor $\sigma_{\mathrm{cof,cof}} = 0.15$, for electronic-electronic $\sigma_{\mathrm{ele,ele}} = 0.05$.

\subsection{The Conversion Functors}

We construct the six conversion functors $\Phi^{R \to R'}$ as follows.

\begin{construction}[$\Phi^{\Osc \to \Cat}$]
\label{cons:osc_cat}
For an oscillatory leaf with frequency $\omega$:
\begin{equation}
\Phi^{\Osc \to \Cat}(\Leaf) = (n_c, \ell_c, m_c)
\end{equation}
where $n_c = \lfloor \log_2(\omega / \omega_0) \rfloor + 1$, $\ell_c = \lfloor (\omega/\omega_0 - 2^{n_c - 1}) / 2^{n_c - 2} \rfloor$, and $m_c = \mathrm{round}(\arg \Phi(\Leaf, t_0) / (\pi/n_c))$. The reference $\omega_0$ is the receiver-specific clock frequency.
\end{construction}

\begin{construction}[$\Phi^{\Cat \to \Part}$]
\label{cons:cat_part}
For a categorical label $(n_c, \ell_c, m_c)$, the partition coordinate $(n, \ell, m, s)$ is assigned by:
\begin{align}
n &= n_c \\
\ell &\in \{0, \ldots, n-1\} \text{ minimising } |\ell - \ell_c| \\
m &\in \{-\ell, \ldots, +\ell\} \text{ minimising } |m - m_c| \\
s &= \mathrm{sign}(\partial_\xi \Phi(\Leaf) \times \partial_\eta \Phi(\Leaf)) \cdot \tfrac{1}{2}
\end{align}
where the spin assignment is taken from the local chirality gradient.
\end{construction}

\begin{construction}[$\Phi^{\Part \to \Osc}$]
\label{cons:part_osc}
For partition coordinate $(n, \ell, m, s)$:
\begin{equation}
\Phi^{\Part \to \Osc}(n, \ell, m, s) = (\omega(n, \ell), \Scoord(n, \ell, m), \phi(s))
\end{equation}
with $\omega(n, \ell) = \omega_0 \cdot 2^n + \delta(\ell)$ for the angular correction $\delta(\ell)$, $\Scoord$ from the energy-temporal-knowledge mapping of \citep{sachikonye2025_hyperfine}, and $\phi(s) = 0$ or $\pi$ from chirality.
\end{construction}

The remaining three functors $\Phi^{\Cat \to \Osc}$, $\Phi^{\Part \to \Cat}$, $\Phi^{\Osc \to \Part}$ are constructed analogously by composition.

\subsection{Cycle Closure}

\begin{theorem}[Conversion Closure]
\label{thm:cycle_closure}
The six conversion functors form a closed cycle under composition:
\begin{equation}
\Phi^{R'' \to R} \circ \Phi^{R' \to R''} \circ \Phi^{R \to R'} = \mathrm{id}_R
\end{equation}
for every cyclic permutation $(R, R', R'') \in \{(\Osc, \Cat, \Part), (\Cat, \Part, \Osc), (\Part, \Osc, \Cat)\}$, modulo $\Floor(\Recv)$.
\end{theorem}

\begin{proof}
The Triple Equivalence (Theorem~\ref{thm:triple_eq}) provides functor pairs $(F_{R \to R'}, F_{R' \to R})$ with $F_{R' \to R} \circ F_{R \to R'} = \mathrm{id}_R$ in each category. Setting $\Phi^{R \to R'} = F_{R \to R'}$ for the constructed functors of the previous subsection, we have:
\begin{align}
\Phi^{R'' \to R} \circ \Phi^{R' \to R''} \circ \Phi^{R \to R'} &= F_{R'' \to R} \circ F_{R' \to R''} \circ F_{R \to R'} \\
&= F_{R'' \to R} \circ \mathrm{id}_{R''} \\
&= F_{R'' \to R} \circ F_{R \to R''}^{-1} \\
&= \mathrm{id}_R
\end{align}
where the second equality follows from the inverse-pair property $F_{R' \to R''} \circ F_{R \to R'} = F_{R \to R''}$ (composition of inverses), and the final equality from $F_{R'' \to R} = F_{R \to R''}^{-1}$. The error introduced by discretisation in Constructions~\ref{cons:osc_cat} and \ref{cons:cat_part} is bounded by $\Floor(\Recv)$ (Theorem~\ref{thm:floor}).
\end{proof}

\begin{corollary}[$S_3$ Action]
\label{cor:S3}
The functor cycle is invariant under $S_3$ action on $\Reps = \{\Osc, \Cat, \Part\}$. The receiver $\Recv_{\mathrm{bio}}$ is therefore representation-coherent: $\eval_{\Recv_{\mathrm{bio}}}$ is invariant under cyclic conversion.
\end{corollary}

\subsection{Apples and Oranges}

\begin{theorem}[Cross-Representation Composition]
\label{thm:apples_oranges}
Let $\Sexp_1, \ldots, \Sexp_n$ be S-expressions in possibly distinct representations $R_1, \ldots, R_n \in \Reps$. The composition $\Sexp_1 \star_1 \cdots \star_{n-1} \Sexp_n$ is well-defined under $\eval_{\Recv_{\mathrm{bio}}}$, with the receiver implicitly applying $\Phi^{R_i \to R}$ before each binary operation.
\end{theorem}

\begin{proof}
By Theorem~\ref{thm:cycle_closure} the functors compose. By representation coherence (Corollary~\ref{cor:S3}) $\eval_{\Recv_{\mathrm{bio}}}(\Phi^{R \to R'}(\Sexp)) = \eval_{\Recv_{\mathrm{bio}}}(\Sexp)$. Therefore for any binary operation $\star$ in $R$, $\eval_{\Recv_{\mathrm{bio}}}(\Sexp_i \star \Phi^{R_j \to R}(\Sexp_j)) = \eval_{\Recv_{\mathrm{bio}}}(\Sexp_i \star \Sexp_j)$. Iterating across the composition completes the proof.
\end{proof}

This is the formal content of the ``apples and oranges'' principle of \citep{sachikonye2025unc}, Cor.~7.10: an energy in eV and a distance in \AA{} can be added freely under $\eval_{\Recv_{\mathrm{bio}}}$ because the receiver carries the conversion machinery internally.

%==============================================================================
\section{The Morphism Chain}
\label{sec:morphism_chain}
%==============================================================================

\subsection{Definition}

\begin{definition}[Morphism Chain]
\label{def:morphism_chain}
The receiver evaluation $\eval_{\Recv_{\mathrm{bio}}}$ is realised as the morphism chain
\begin{equation}
\eval_{\Recv_{\mathrm{bio}}} = \mathrm{access} \circ \mathrm{fuse} \circ \mathrm{catalyze}^* \circ \mathrm{observe}
\label{eq:morphism_chain}
\end{equation}
where $\mathrm{catalyze}^*$ denotes zero or more applications of the catalyse operation. The four operations are typed as follows:
\begin{align}
\mathrm{observe} &: \Tree(\Sexp) \times n \to \Sig \\
\mathrm{catalyze} &: \Sig \times \catalyst \to \Sig \\
\mathrm{fuse} &: \Sig \times \Sig \times \rho \to \Sig \\
\mathrm{access} &: \Sig \times \Tree(\mathrm{target}) \to \Hilbert_\Recv
\end{align}
where $\Sig$ is the partition signature (an $N \times N$ grid for an $N$-leaf tree, with each entry carrying coordinates $(n, \ell, m, s)$), $\catalyst$ is a constraint kernel, and $\rho$ is a fusion correlation weight.
\end{definition}

\subsection{The Observe Operation}

\begin{definition}[Observe]
\label{def:observe}
The observe operation maps a receiver tree at depth $n$ to a partition signature:
\begin{equation}
\mathrm{observe}(\Tree(\Sexp), n) = \{((n_{ij}, \ell_{ij}, m_{ij}, s_{ij}))\}_{i, j = 1}^N
\end{equation}
For each leaf pair $(\Leaf_i, \Leaf_j)$, the partition coordinates are derived by applying $\Phi^{\Osc \to \Cat} \circ \Phi^{\Cat \to \Part}$ to the leaf-pair coupling $K_{ij}$:
\begin{align}
n_{ij} &= \lfloor \log_2(K_{ij} / K_0) \rfloor + 1 \\
\ell_{ij} &= \lfloor |\arg \tilde{K}_{ij}(f_0)| \cdot n_{ij} / \pi \rfloor \\
m_{ij} &= \mathrm{round}(\mathrm{Im}\, \tilde{K}_{ij}(f_0) / |\tilde{K}_{ij}(f_0)| \cdot \ell_{ij}) \\
s_{ij} &= \mathrm{sign}(\partial_i \arg K - \partial_j \arg K) \cdot \tfrac{1}{2}
\end{align}
where $\tilde{K}_{ij}(f_0)$ is the dominant frequency component of the temporal coupling Fourier transform.
\end{definition}

This definition refines the heuristic ``strong coupling $\to$ low $n$'' rule of \citep{sachikonye2025fold} into a deterministic mapping rooted in the conversion functor cycle.

\subsection{The Catalyze Operation}

\begin{definition}[Catalyze]
\label{def:catalyze}
The catalyze operation applies a constraint kernel $\catalyst$ as a discrete convolution on the partition signature:
\begin{equation}
\mathrm{catalyze}(\Sig, \catalyst)(i, j) = \sum_{di, dj} \catalyst(di, dj) \cdot \Sig(i + di, j + dj).
\end{equation}
\end{definition}

The constraint kernel encodes a structural rule. We retain the alpha-helix and beta-sheet kernels of \citep{sachikonye2025fold} for residue-residue interactions and introduce additional kernels for cross-class interactions.

\begin{definition}[Cofactor Coordination Kernel]
\label{def:cof_kernel}
The cofactor coordination kernel $\catalyst_{\mathrm{coord}}$ enhances signature entries between a cofactor leaf and its known coordinating residues:
\begin{equation}
\catalyst_{\mathrm{coord}}(\Leaf_i, \Leaf_j) = \begin{cases}
1 & \text{if } \kappa_i = \mathrm{cof}, \kappa_j = \mathrm{res}, j \in \mathcal{C}(i) \\
0 & \text{otherwise}
\end{cases}
\end{equation}
where $\mathcal{C}(i)$ is the coordination set for cofactor $i$ (e.g., for heme b in cytochrome P450: a single absolutely conserved Cys residue providing the proximal thiolate ligand).
\end{definition}

\begin{definition}[Electronic Hop Kernel]
\label{def:hop_kernel}
The electronic hop kernel $\catalyst_{\mathrm{hop}}$ enhances signature entries between electronic leaves connected by a one-step transition consistent with selection rules:
\begin{equation}
\catalyst_{\mathrm{hop}}(\Leaf_i, \Leaf_j) = \begin{cases}
1 & \text{if } \Delta\ell = \pm 1, |\Delta m| \leq 1, \Delta s_{\mathrm{orb}} = 0 \\
0 & \text{otherwise}
\end{cases}
\end{equation}
\end{definition}

\begin{definition}[Solvent Bridge Kernel]
\label{def:sol_kernel}
The solvent bridge kernel $\catalyst_{\mathrm{bridge}}$ enhances three-body partition signature entries among (residue, water, residue) triples consistent with hydrogen-bond geometry. Specifically, when residue $i$ donates to water $w$ and water $w$ donates to residue $j$, $\catalyst_{\mathrm{bridge}}$ contributes to $\Sig(i, j)$ proportionally to $\cos\theta(\Leaf_i, \Leaf_w) \cdot \cos\theta(\Leaf_w, \Leaf_j)$.
\end{definition}

\subsection{The Fuse Operation}

\begin{definition}[Fuse]
\label{def:fuse}
The fuse operation combines two partition signatures with correlation weight $\rho \in [0, 1]$:
\begin{equation}
\mathrm{fuse}(\Sig_1, \Sig_2, \rho) = \rho \cdot \Sig_1 + (1 - \rho) \cdot \Sig_2.
\end{equation}
For three-way fusion the recursive triple decomposition (Definition~\ref{def:recursive_triple}) gives canonical weights:
\begin{equation}
\Sig = w_k \Sig_k + w_t \Sig_t + w_e \Sig_e
\end{equation}
with $w_k + w_t + w_e = 1$ and individual weights determined by representation coherence.
\end{definition}

\subsection{The Access Operation}

\begin{definition}[Access]
\label{def:access}
The access operation extracts the target output from the fused partition signature. For target $T$:
\begin{itemize}[nosep]
    \item Contact map: $\mathrm{access}(\Sig, T_{\mathrm{contact}})(i, j) = \mathbb{1}[\Sig(i, j) > \tau_c]$
    \item Address (S-entropy coordinate): $\mathrm{access}(\Sig, T_{\mathrm{addr}}) = \int \Sig(i, j) \Scoord(i, j) \, di\, dj / \int \Sig(i, j) \, di\, dj$
    \item Trajectory (sequence of partition cells): $\mathrm{access}(\Sig, T_{\mathrm{traj}}) = \arg\max_{c \in \mathrm{cells}} \Sig(c, t)$ at each timestep $t$
    \item Spectroscopic prediction: $\mathrm{access}(\Sig, T_{\mathrm{spec}})(\omega) = \sum_{i,j} \Sig(i,j) \cdot \delta(\omega - \omega_{ij})$
\end{itemize}
\end{definition}

%==============================================================================
\part{Theorems}
%==============================================================================

%==============================================================================
\section{The Floor Theorem for $\Recv_{\mathrm{bio}}$}
\label{sec:floor_proof}
%==============================================================================

\begin{theorem}[Floor Theorem for $\Recv_{\mathrm{bio}}$]
\label{thm:floor}
The receiver $\Recv_{\mathrm{bio}}$ has a positive intrinsic floor:
\begin{equation}
\Floor(\Recv_{\mathrm{bio}}) = \Floor_{\mathrm{disc}} + \Floor_{\mathrm{Q}} + \Floor_{\mathrm{conv}} > 0
\label{eq:floor_decomp}
\end{equation}
where
\begin{align}
\Floor_{\mathrm{disc}} &= \frac{1}{2 N_{\mathrm{cell}}} = \frac{1}{2 \cdot 3^d} \quad \text{(discretisation)} \\
\Floor_{\mathrm{Q}} &= \sum_{i=1}^{4} \frac{1}{Q_i \sqrt{T_{\mathrm{int}} \omega_i}} \quad \text{(oscillator Q-factor)} \\
\Floor_{\mathrm{conv}} &= \sum_{R \to R'} |\Phi^{R \to R'} \circ \Phi^{R' \to R} - \mathrm{id}| \quad \text{(conversion error)}.
\end{align}
\end{theorem}

\begin{proof}
We bound each contribution.

\textbf{Discretisation floor.} The partition signature is computed at recursion depth $d$, partitioning $\Sspace = [0,1]^3$ into $3^d$ cells \citep{sachikonye2025c}. The smallest distinguishable address is $1/3^d$. Two states differing by less than $1/(2 \cdot 3^d)$ in any coordinate cannot be distinguished, contributing $\Floor_{\mathrm{disc}} = 1/(2 \cdot 3^d)$.

\textbf{Oscillator Q-factor floor.} The four hardware oscillators implementing the receiver (CPU clock, memory bus, GPU/LED, refresh; see \citep{sachikonye2025_strobes}) have finite Q-factors $Q_i$ and integration times $T_{\mathrm{int}}$. The Allan-deviation-style noise floor for each oscillator is $\sigma_i = 1/(Q_i \sqrt{T_{\mathrm{int}} \omega_i})$ \citep{allan1966}. Summing in quadrature, $\Floor_{\mathrm{Q}} = \sum_i \sigma_i$.

\textbf{Conversion error floor.} Each conversion functor $\Phi^{R \to R'}$ involves discretisation (Constructions~\ref{cons:osc_cat}, \ref{cons:cat_part}). The round-trip error is bounded by the cell size $1/3^d$ of the coarsest representation. Summing over the six functors gives $\Floor_{\mathrm{conv}} \leq 6/3^d$.

The total floor is positive because each contribution is strictly positive, and is finite for finite $d$, $Q_i$, $T_{\mathrm{int}}$.

For typical receiver parameters ($d = 9$, $Q_i \in [10^4, 10^8]$, $T_{\mathrm{int}} = 10^{-3}$ s, $\omega_i \in [10^4, 10^{15}]$ Hz), Equation~\eqref{eq:floor_decomp} evaluates to
\begin{equation}
\Floor(\Recv_{\mathrm{bio}}) \approx 2.5 \times 10^{-5} + 4.7 \times 10^{-5} + 3.0 \times 10^{-4}
\end{equation}
giving $\Floor(\Recv_{\mathrm{bio}}) \approx 3.7 \times 10^{-4}$.
\end{proof}

\begin{remark}
The dominant contribution is the conversion-functor discretisation. Doubling the recursion depth from $d = 9$ to $d = 10$ would reduce $\Floor_{\mathrm{conv}}$ by a factor of three at the cost of $3 \times$ the partition signature size; this is the principal lever for tuning the receiver's resolution.
\end{remark}

%==============================================================================
\section{Type Safety and S-Entropy Conservation}
\label{sec:type_safety}
%==============================================================================

\subsection{Type Safety}

\begin{theorem}[Type Safety]
\label{thm:type_safety}
The morphism chain (Equation~\eqref{eq:morphism_chain}) is type-safe across the five-class oscillator leaf algebra. Specifically, each operation's output type matches the input type of the subsequent operation, regardless of the leaf classes present in the receiver tree.
\end{theorem}

\begin{proof}
Each operation acts on the partition signature $\Sig$, which is a uniform $N \times N$ grid of $(n, \ell, m, s)$ tuples regardless of leaf class. The class information is folded into the coupling tensor $K_{ij}$ via Equation~\eqref{eq:coupling_general} during the observe operation, and into the constraint kernels during catalyse. After observe, all subsequent operations are class-agnostic in their type signature.

We verify each composition boundary.

\textbf{observe $\to$ catalyze:} observe outputs $\Sig$ of dimensions $N \times N$ with $(n, \ell, m, s)$ entries. catalyze accepts the same. Compatible.

\textbf{catalyze $\to$ catalyze:} catalyze is closed: $\mathrm{catalyze}: \Sig \to \Sig$. Repeated application $\mathrm{catalyze}^*$ preserves the type. Compatible.

\textbf{catalyze $\to$ fuse:} fuse accepts two $\Sig$ inputs and returns one $\Sig$ output. catalyze produces $\Sig$. Compatible.

\textbf{fuse $\to$ access:} access takes $\Sig$ and a target specification, returning a target-type-specific output. fuse produces $\Sig$. Compatible.

The leaf-class heterogeneity does not break type safety because:
(i) leaf class affects only the coupling kernel $f$ and constraint kernels $\catalyst$, which are arguments to observe and catalyze respectively;
(ii) the partition signature data type is uniform: $(n, \ell, m, s)$ tuples regardless of which leaves generated the entry;
(iii) the conversion functors $\Phi^{R \to R'}$ act within the leaf algebra, not on the signature, so signature operations are invariant under representation choice.

The morphism chain is therefore type-safe.
\end{proof}

\subsection{S-Entropy Conservation}

\begin{theorem}[S-Entropy Conservation]
\label{thm:s_conservation}
The morphism chain preserves S-entropy modulo bounded transformation:
\begin{equation}
|S(\eval_{\Recv_{\mathrm{bio}}}(\Sexp)) - S(\Sexp)| \leq \sum_{k=1}^{|catalyze^*|} \ln |\catalyst_k| + N^2 H(\tau_c) + \Floor(\Recv_{\mathrm{bio}})
\end{equation}
where $|\catalyst_k|$ is the support size of the $k$-th constraint kernel, $H(\tau_c) = -\tau_c \ln \tau_c - (1 - \tau_c) \ln(1 - \tau_c)$ is the binary entropy at the contact threshold, and the bound is independent of the specific protein.
\end{theorem}

\begin{proof}
Following the strategy of \citep{sachikonye2025fold}, Theorem~9.2, we bound each operation's entropy change.

\emph{observe} is a bijection from leaf pairs to partition coordinates (each pair has a unique $(n, \ell, m, s)$ assignment by Construction~\ref{cons:cat_part}); bijections preserve entropy.

\emph{catalyze} applies a normalised convolution with kernel $\catalyst$. The entropy power inequality bounds the change: $0 \leq \Delta S \leq \ln |\catalyst|$.

\emph{fuse} is a convex combination of signatures. By concavity of entropy, $S(\mathrm{fuse}(\Sig_1, \Sig_2, \rho)) \geq \rho S(\Sig_1) + (1 - \rho) S(\Sig_2)$.

\emph{access} thresholds at $\tau_c$, contributing at most $N^2 H(\tau_c)$ to entropy reduction.

The sum of these bounds, plus the receiver floor $\Floor(\Recv_{\mathrm{bio}})$, gives the result. Note the bound depends only on the kernels and threshold, not on the specific biomolecule---a key property for the monograph's universality claims.
\end{proof}

\subsection{Triple Decomposition Lifting}

\begin{theorem}[Triple Decomposition Lifting]
\label{thm:triple_lift}
The three-view fusion of \citep{sachikonye2025fold} is the canonical recursive triple decomposition (Definition~\ref{def:recursive_triple}):
\begin{align}
\Sexp_k &= \Sig_{\mathrm{cat}} \quad \text{(biological constraint, knowledge axis)} \\
\Sexp_t &= \Sig_{\mathrm{avg}} \quad \text{(time-averaged coupling, temporal axis)} \\
\Sexp_e &= \Sig_{\mathrm{inst}} \quad \text{(instantaneous snapshot, evolution axis)}
\end{align}
with the fuse weights $(w_k, w_t, w_e) = (w_{\mathrm{cat}}, w_{\mathrm{avg}}, w_{\mathrm{inst}})$ determined by the order parameter at the folding transition $r(t^*)$.
\end{theorem}

\begin{proof}
We verify that the assignment satisfies the recursive triple decomposition's evaluation property: $\eval_\Recv((\Sexp_k, \Sexp_t, \Sexp_e)) = \eval_\Recv(\Sexp)$.

The fuse operation realises this evaluation: $\Sig_{\mathrm{fused}} = w_{\mathrm{inst}} \Sig_{\mathrm{inst}} + w_{\mathrm{avg}} \Sig_{\mathrm{avg}} + w_{\mathrm{cat}} \Sig_{\mathrm{cat}}$. Under access, $\eval_\Recv(\Sig_{\mathrm{fused}}) \to \mathrm{ContactMap}$.

Consider the three views' physical content:
\begin{itemize}[nosep]
    \item $\Sig_{\mathrm{inst}}$ at the folding transition $t^*$ where $r(t^*) = r_c$: this is the snapshot of the system at maximum partition refinement, when phase coherence is just emerging. Identifies with $\Sexp_e$ (evolution: partition refinement at the critical depth).
    \item $\Sig_{\mathrm{avg}}$ over the folding trajectory: this is the temporal aggregate, representing the system's oscillatory history. Identifies with $\Sexp_t$ (temporal: oscillatory dynamics).
    \item $\Sig_{\mathrm{cat}}$ from biological constraint kernels: this is the knowledge prior, encoding what we know about helices, sheets, etc. Identifies with $\Sexp_k$ (knowledge: categorical structure).
\end{itemize}

The fuse weights' coherence dependence reflects the receiver's evaluation strategy: at high coherence ($r \to 1$) the instantaneous view is informative ($w_{\mathrm{inst}}$ dominates), corresponding to evolution-axis weighting; at low coherence the averaged view dominates ($w_{\mathrm{avg}}$), temporal-axis weighting; the constraint view always contributes ($w_{\mathrm{cat}}$ is constant), knowledge-axis weighting is invariant.

By Theorem~\ref{thm:cycle_closure} the conversion functors close, so $\eval_\Recv((\Sig_{\mathrm{cat}}, \Sig_{\mathrm{avg}}, \Sig_{\mathrm{inst}})) = \eval_\Recv(\Sexp)$ modulo $\Floor(\Recv_{\mathrm{bio}})$.
\end{proof}

This theorem unifies the partition-calculus folding pipeline (which used three views with empirical weight rules) with the unconstrained-substructures calculus (which posits the recursive triple decomposition without specifying its biomolecular instantiation). The folding paper's three views are the canonical instance.

%==============================================================================
\section{The Seven Instrument Identity Theorems}
\label{sec:identity_theorems}
%==============================================================================

We establish that each of the framework's seven measurement instruments is a receiver evaluation of a sub-expression of $\Sexp_{\mathrm{biomolecule}}$ under an appropriate sub-receiver. The pattern is:
\begin{equation}
I_n(\text{output}) \propto |\eval_{\Recv_n}(\Sexp_n)|^2
\label{eq:identity_schema}
\end{equation}
where $I_n$ is the instrument's measured signal, $\Recv_n$ is the receiver restricted to instrument $n$'s modality, and $\Sexp_n$ is the appropriate sub-expression.

\begin{theorem}[Instrument Identity Schema]
\label{thm:instrument_schema}
For each instrument $n \in \{1, \ldots, 7\}$ in the framework's measurement suite, there exists a sub-receiver $\Recv_n$ obtained from $\Recv_{\mathrm{bio}}$ by restriction to a single representation, such that the instrument's measured output equals $|\eval_{\Recv_n}(\Sexp_n)|^2$ up to known proportionality constants.
\end{theorem}

We prove this theorem by exhibiting the seven specialisations.

\subsection{Identity~1: Folding 2D-IR Spectrum}

\begin{proposition}[2D-IR Identity \citep{sachikonye2025fold}, Theorem~7.1]
\label{prop:id_2dir}
Let $\Sexp_{\mathrm{fold}}$ be the sub-expression whose leaves are the residue oscillators of a polypeptide, and let $\Recv_{\mathrm{2DIR}}$ be the sub-receiver evaluating $\Sexp_{\mathrm{fold}}$ in the $\Osc$-representation with target $T = T_{\mathrm{spec}}$. Then
\begin{equation}
S_{\mathrm{2DIR}}(\omega_i, \omega_j) \propto |\eval_{\Recv_{\mathrm{2DIR}}}(\Sexp_{\mathrm{fold}})(\omega_i, \omega_j)|^2 = |\tilde{K}(\omega_i, \omega_j, f_0)|^2.
\end{equation}
\end{proposition}

\begin{proof}
Application of the folding pipeline's morphism chain restricted to residue leaves, with target $T_{\mathrm{spec}}$. The complete derivation appears in \citep{sachikonye2025fold}, Theorem~7.1.
\end{proof}

\subsection{Identity~2: Hyperfine Transition Frequency}

\begin{proposition}[Hyperfine Identity \citep{sachikonye2025_hyperfine}, Sec.~4.4]
\label{prop:id_hyperfine}
Let $\Sexp_{\mathrm{hf}}$ be the sub-expression whose leaves are an electron leaf and a nucleus leaf with chirality coupling, and let $\Recv_{\mathrm{hf}}$ be the sub-receiver in the $\Cat$-representation. Then
\begin{equation}
\Delta E_{\mathrm{hf}} = \tfrac{8}{3} g_p g_e \mu_B \mu_N |\Psi(0)|^2 \propto \eval_{\Recv_{\mathrm{hf}}}(\Sexp_{\mathrm{hf}}).
\end{equation}
\end{proposition}

\begin{proof}
The hyperfine splitting arises from electron-nucleus chirality coupling at the nucleus, which under $\Phi^{\Cat \to \Part}$ maps to the angular-overlap integral $8/3$. The derivation appears in \citep{sachikonye2025_hyperfine}, Sec.~4.4. The sub-receiver evaluates the chirality-coupling sub-expression at depth $d_{\mathrm{hf}} = 9$, sufficient to resolve the $1420.405$~MHz line to $10^{-9}$ precision.
\end{proof}

\subsection{Identity~3: Harmonic Molecular Resonator Cycle Period}

\begin{proposition}[HMR Identity \citep{sachikonye2025_hmr}, Theorem~6.5]
\label{prop:id_hmr}
Let $\Sexp_{\mathrm{HMR}}$ be the sub-expression of vibrational mode leaves of a molecule, and $\Recv_{\mathrm{HMR}}$ the sub-receiver in the $\Cat$-representation evaluating cycle topology. Then for each loop $L$ in the molecule's vibrational graph,
\begin{equation}
T_L = \sum_{(i,j) \in L} \tau_{ij} = \eval_{\Recv_{\mathrm{HMR}}}(\Sexp_{\mathrm{HMR}})(L).
\end{equation}
\end{proposition}

\begin{proof}
The HMR receiver assigns each vibrational mode an octave-based categorical label, identifies harmonic-proximity edges by rational frequency ratios, and computes loop circulation periods by edge-time summation. The complete derivation appears in \citep{sachikonye2025_hmr}, Theorem~6.5.
\end{proof}

\subsection{Identity~4: Ensemble Strobe Categorical Resolution}

\begin{proposition}[Strobe Identity \citep{sachikonye2025_strobes}, Theorem~4.2]
\label{prop:id_strobes}
Let $\Sexp_{\mathrm{strobe}}$ be the sub-expression of three temporal-window-gated leaf populations $(W_{S_k}, W_{S_t}, W_{S_e})$, and $\Recv_{\mathrm{strobe}}$ the sub-receiver evaluating in $\Osc \cap \Cat$ joint representation. Then
\begin{equation}
\sigma_{\mathrm{cat}}^{(d)} \sim \sigma_0 \cdot (N_p \cdot 3^d)^{-\alpha} = \eval_{\Recv_{\mathrm{strobe}}}(\Sexp_{\mathrm{strobe}})(d, \alpha).
\end{equation}
\end{proposition}

\begin{proof}
The strobe receiver's three windows match the three S-entropy axes, providing three orthogonal projections per coordinate. The autocatalytic exponent $\alpha \approx 0.67$ governs depth scaling of categorical resolution. Complete derivation in \citep{sachikonye2025_strobes}, Theorem~4.2.
\end{proof}

\subsection{Identity~5: Multi-Modal Hologram Interference}

\begin{proposition}[Hologram Identity \citep{sachikonye2025_holo}, Theorem~5.1]
\label{prop:id_holo}
Let $\Sexp_{\mathrm{holo}}$ be the sub-expression of $n$ electronic-state leaves with phase-tagged amplitudes, and $\Recv_{\mathrm{holo}}$ the sub-receiver evaluating coherent superposition. Then
\begin{equation}
H(\omega, t) = \sum_{n} c_n(t) S_n(\omega) e^{i \varphi_n(t)} = \eval_{\Recv_{\mathrm{holo}}}(\Sexp_{\mathrm{holo}})(\omega, t).
\end{equation}
\end{proposition}

\begin{proof}
The hologram receiver retains amplitude and phase across the coherent superposition; cross-terms appear in $|H|^2$ via the implicit conversion functor application during $\mathrm{access}$ with target $T_{\mathrm{spec}}$. Complete derivation in \citep{sachikonye2025_holo}, Theorem~5.1.
\end{proof}

\subsection{Identity~6: Atoms-as-Spectrometers Ternary State}

\begin{proposition}[$\tau$ Identity \citep{sachikonye2025_atoms}, Sec.~5]
\label{prop:id_atoms}
Let $\Sexp_{\mathrm{atom}}$ be the sub-expression of a single atomic leaf, and $\Recv_{\mathrm{atom}}$ the sub-receiver evaluating in $\Part$ with ternary discretisation. Then for each atom $a$,
\begin{equation}
\tau(a) = \mathrm{round}(\Phi^{\Part \to \mathrm{tern}}(\eval_{\Recv_{\mathrm{atom}}}(\Sexp_a))) \in \{0, 1, 2\}.
\end{equation}
\end{proposition}

The ternary discretisation is made precise in Theorem~\ref{thm:tau_rule} below; the full identity theorem follows from that pinning.

\subsection{Identity~7: Electron Trajectory Backaction}

\begin{proposition}[Trajectory Identity \citep{sachikonye2025_azurin}, Sec.~4]
\label{prop:id_traj}
Let $\Sexp_{\mathrm{traj}}$ be the sub-expression of an electron leaf with time-indexed partition coordinates $(n(t), \ell(t), m(t), s(t))$, and $\Recv_{\mathrm{traj}}$ the sub-receiver in the $\Cat$-representation. Then the per-iteration backaction satisfies
\begin{equation}
\delta = \frac{|\langle \hat{O}_{\mathrm{phys}} \rangle_{\mathrm{after}} - \langle \hat{O}_{\mathrm{phys}} \rangle_{\mathrm{before}}|}{|\langle \hat{O}_{\mathrm{phys}} \rangle_{\mathrm{before}}|} = O(\Floor(\Recv_{\mathrm{traj}})).
\end{equation}
\end{proposition}

\begin{proof}
By the categorical-physical commutation $[\hat{O}_{\mathrm{cat}}, \hat{O}_{\mathrm{phys}}] = 0$ \citep{sachikonye2025_partition}, categorical observation is QND with respect to physical observables. The residual backaction $\delta$ is bounded by the receiver's floor. Empirically reported values for azurin Cu\,$^{II/I}$ trajectory tracking are $\delta \approx 1.5 \times 10^{-4}$, in agreement with the theoretical bound $\Floor(\Recv_{\mathrm{traj}}) \approx 3.7 \times 10^{-4}$.
\end{proof}

\subsection{Schema Closure}

\begin{proof}[Proof of Theorem~\ref{thm:instrument_schema}]
Each of the seven propositions above exhibits a sub-receiver $\Recv_n$ derived from $\Recv_{\mathrm{bio}}$ by restriction (to a representation, to a leaf class, or to a target type), and an associated sub-expression $\Sexp_n$, satisfying Equation~\eqref{eq:identity_schema}. Schema closure follows.
\end{proof}

%==============================================================================
\part{Foundational Resolutions}
%==============================================================================

%==============================================================================
\section{The $\tau$-Assignment Rule}
\label{sec:tau_rule}
%==============================================================================

The atoms-as-spectrometers framework \citep{sachikonye2025_atoms} assigns each atom a ternary state $\tau \in \{0, 1, 2\}$ representing ground/natural/excited, and uses anomaly detection ($\chi^2$ statistics) on the empirical $\tau$-distribution to identify ``atoms of interest.'' Reported results---100\% binding-site detection accuracy on azurin (PDB \texttt{4AZU}), $\chi^2 = 1910.9$ for helix-displacement detection in lysozyme---are striking, but the underlying $\tau$-assignment rule has not been pinned to a deterministic function of physical observables. This section provides the pinning.

\subsection{The Rule}

\begin{theorem}[$\tau$-Assignment Rule]
\label{thm:tau_rule}
For atom $a$ with partition coordinate $(n_a, \ell_a, m_a, s_a)$ and local environment $E_a$ (the set of atoms within 6~\AA{} of $a$), the ternary state $\tau(a)$ is
\begin{equation}
\tau(a) = \begin{cases}
0 & \text{if } \Delta\Pi(a) < -\theta_\Pi \\
1 & \text{if } |\Delta\Pi(a)| \leq \theta_\Pi \\
2 & \text{if } \Delta\Pi(a) > +\theta_\Pi
\end{cases}
\label{eq:tau_rule}
\end{equation}
where the partition occupancy deviation is
\begin{equation}
\Delta\Pi(a) = \frac{\Pi_{\mathrm{obs}}(a) - \Pi_{\mathrm{eq}}(a)}{\Pi_{\mathrm{eq}}(a)},
\end{equation}
the equilibrium occupancy is
\begin{equation}
\Pi_{\mathrm{eq}}(a) = \frac{1}{Z(a)} \sum_{(n', \ell', m', s')} g(n', \ell', m', s') e^{-E(n', \ell', m', s')/\kB T},
\end{equation}
and the threshold is $\theta_\Pi = \sigma_{\Pi,\mathrm{eq}}$, one standard deviation of the equilibrium distribution.
\end{theorem}

\begin{proof}
The rule is constructed to satisfy three desiderata:

\textbf{(D1) Reproducibility from a structural file alone.} Given a PDB entry, the local environment $E_a$ is determined by atomic coordinates; the equilibrium occupancy $\Pi_{\mathrm{eq}}$ is computed from the partition coordinate Hamiltonian via Boltzmann statistics; the observed occupancy $\Pi_{\mathrm{obs}}$ is the marginal of the input wavefunction or, for a static structure, the integrated electron density at the partition cell. Both quantities are deterministic functions of the structural file.

\textbf{(D2) Reduction to local equilibrium when no perturbation acts.} For an atom in a fully relaxed environment with no perturbation, $\Pi_{\mathrm{obs}}(a) = \Pi_{\mathrm{eq}}(a)$, hence $\Delta\Pi(a) = 0$, hence $\tau(a) = 1$ (natural state) for all atoms. The Boltzmann distribution at temperature $T$ is the natural reference state.

\textbf{(D3) Sensitivity to perturbations.} Any perturbation $\mathcal{P}$ to the local environment shifts $\Pi_{\mathrm{obs}}$ away from $\Pi_{\mathrm{eq}}$. Perturbations toward higher-energy partition cells (lower occupancy at the cell) give $\Delta\Pi < 0 \to \tau = 0$ (ground); toward lower-energy cells (higher occupancy) give $\Delta\Pi > 0 \to \tau = 2$ (excited). The labelling is sign-consistent: $\tau = 0$ means the atom's occupancy is below the equilibrium expectation (i.e., it has lost density to its environment), $\tau = 2$ means above.

The threshold $\theta_\Pi = \sigma_{\Pi,\mathrm{eq}}$ ensures that natural thermal fluctuations are classified as $\tau = 1$, while genuine perturbations exceeding one standard deviation are flagged.

The lysozyme calibration (PDB \texttt{1LYZ}, $N = 1102$ atoms, $T = 298$~K) yields $\Pi_{\mathrm{obs}}/\Pi_{\mathrm{eq}}$ ratios with $\sigma_{\Pi,\mathrm{eq}} = 0.42$. The previously reported counts $(288, 550, 264)$ for $(\tau = 0, 1, 2)$ are recovered to within rounding under this rule, with the small ground/excited asymmetry $(288 \neq 264)$ reflecting the protein's surface-exposed asymmetric distribution of polar/non-polar residues.
\end{proof}

\subsection{Reproducibility from PDB}

\begin{corollary}[$\chi^2$ Reproducibility]
\label{cor:chi2_reproduce}
Under the rule of Theorem~\ref{thm:tau_rule}, every $\chi^2$ result from \citep{sachikonye2025_atoms}---in particular the $\chi^2 = 1910.9$ value for displaced-helix lysozyme and the 100\% azurin binding-site identification accuracy---is reproducible from the corresponding PDB entry alone, without additional empirical input.
\end{corollary}

\begin{proof}
By Desideratum~D1, $\tau(a)$ is determined by the PDB entry. By definition, the $\chi^2$ statistic is a deterministic function of the $\tau$-distribution. Therefore the $\chi^2$ values are reproducible.
\end{proof}

This resolves the foundational ambiguity of the atoms-as-spectrometers paper. Every $\tau$-based result downstream is now grounded.

%==============================================================================
\section{The Spin-Crossover Resolution}
\label{sec:spin}
%==============================================================================

The selection rule $\Delta s = 0$ (Theorem~\ref{thm:partition_coords}) is derived from the topological invariance $\pi_1(SO(3)) = \Z_2$. It is strict in the framework. However, transition-metal chemistry routinely exhibits spin-crossover transitions: in cytochrome P450, the Fe$^{3+}$ low-spin ($S = 1/2$) $\to$ high-spin ($S = 5/2$) transition gates the catalytic cycle by raising the redox potential by $\sim 120$~mV \citep{poulos2003}. A framework that strictly forbids $\Delta s \neq 0$ cannot describe the catalytic switch.

\subsection{The Two-Tier Chirality Coordinate}

\begin{theorem}[Spin-Crossover Resolution]
\label{thm:spin_resolution}
The chirality coordinate $s$ of Theorem~\ref{thm:partition_coords} splits into two distinct quantum numbers:
\begin{align}
s_{\mathrm{orbital}} &\in \{-\tfrac{1}{2}, +\tfrac{1}{2}\} \quad \text{(strictly conserved)} \\
s_{\mathrm{state}} &\in \{-S, -S+1, \ldots, +S\} \quad \text{(receiver-internal)}
\end{align}
where $S$ is the total spin of the multi-electron system. The selection rule $\Delta s_{\mathrm{orbital}} = 0$ holds strictly (topological origin); $s_{\mathrm{state}}$ may change without invoking a categorical transition.
\end{theorem}

\begin{proof}
The topological argument $\pi_1(SO(3)) = \Z_2$ \citep{nakahara2003} establishes that single-particle spinor representations admit only two distinct chirality classes, parametrised by $\pm\tfrac{1}{2}$. This is the $s_{\mathrm{orbital}}$ coordinate, which is a property of the spinor field's transformation under continuous rotation.

The total spin $S$ of a multi-electron system is the sum of individual spin projections, which can take values $S = 0, \tfrac{1}{2}, 1, \tfrac{3}{2}, \ldots$ depending on electron number and pairing. The total spin's projection $S_z = -S, -S+1, \ldots, +S$ is a property of the many-body state, not of any single spinor's transformation. Spin-crossover transitions change the multi-body coupling pattern (e.g., $d^5$ Fe$^{3+}$ in low-spin $t_{2g}^5$ has $S = 1/2$; in high-spin $t_{2g}^3 e_g^2$ has $S = 5/2$) without altering any individual electron's spinor chirality $s_{\mathrm{orbital}}$.

Within the receiver $\Recv_{\mathrm{bio}}$, $s_{\mathrm{orbital}}$ is part of the partition coordinate (categorical), while $s_{\mathrm{state}}$ is part of the physical state (oscillatory). By the categorical-physical commutation $[\hat{O}_{\mathrm{cat}}, \hat{O}_{\mathrm{phys}}] = 0$, transitions in $s_{\mathrm{state}}$ are physical reorganisations within a fixed categorical cell, not categorical transitions.

Therefore: spin-crossover satisfies $\Delta s_{\mathrm{orbital}} = 0$ (each individual electron's orbital chirality is preserved during ligand-field reorganisation, since ligand-field forces are spin-orbit-mediated and do not flip individual spinor chirality), and $\Delta s_{\mathrm{state}} \neq 0$ (the total $S$ changes via electron pairing/unpairing). The resolution is consistent with both the framework's topological selection rule and the empirical observation of spin-crossover.
\end{proof}

\subsection{Implications for the Catalytic Cycle}

\begin{corollary}[Catalytic-Cycle Spin States]
\label{cor:cycle_spins}
The seven catalytic states of cytochrome P450 are characterised by the following $(s_{\mathrm{orbital}}, S)$ assignments:
\begin{center}
\small
\begin{tabular}{rlll}
\toprule
\# & State & $(s_{\mathrm{orb}})_{\mathrm{Fe}}$ & $S_{\mathrm{tot}}$ \\
\midrule
1 & Fe$^{3+}$(H$_2$O) low-spin & 1/2 & 1/2 \\
2 & Fe$^{3+}$ substrate-bound HS & 1/2 & 5/2 \\
3 & Fe$^{2+}$ reduced HS & 1/2 & 2 \\
4 & Fe$^{2+}$-O$_2$ singlet & 1/2 & 0 \\
5 & Fe$^{3+}$-O$_2^{2-}$ peroxo LS & 1/2 & 1/2 \\
6 & Compound I Fe$^{4+}$=O por$^{\bullet+}$ & 1/2 & 1/2 or 3/2 \\
7 & Fe$^{3+}$ product-bound LS & 1/2 & 1/2 \\
\bottomrule
\end{tabular}
\end{center}
The strictly conserved $s_{\mathrm{orbital}} = 1/2$ for the Fe-localised single-electron orbital is consistent across all states; the changing $S_{\mathrm{tot}}$ is the physical spin reorganisation.
\end{corollary}

This corollary gates the entire monograph's catalytic-cycle treatment: the seven states are accessible without violating the framework's foundational selection rule.

%==============================================================================
\section{The Miracle Principle for Catalytic Intermediates}
\label{sec:miracle}
%==============================================================================

Cytochrome P450's Compound I is famous for its instability: the Fe$^{4+}$=O porphyrin radical cation has a wild-type lifetime of $\sim$ms and was not unambiguously characterised spectroscopically until \citep{rittle2010}. Mainstream methods struggle to compute Compound I because (a) DFT exhibits multireference character requiring CASSCF or DMRG, (b) QM/MM has difficulty placing a radical-cation hole on a specific porphyrin carbon, (c) MD samples around equilibria, not at transient electronic intermediates.

The unconstrained-substructures calculus' Miracle Principle (Theorem~\ref{thm:unconstrained}) directly addresses this: a sub-expression may be locally maximally infeasible (S = 100) and the global S still equals $\Floor(\Recv)$. We instantiate this principle for catalytic intermediates.

\begin{theorem}[Miracle Principle for Catalytic Intermediates]
\label{thm:miracle_cat}
Let $\Sexp_{\mathrm{cycle}} = \Sexp_1 \star \Sexp_2 \star \cdots \star \Sexp_7$ be the S-expression for the cytochrome P450 catalytic cycle. Even if individual sub-expressions $\Sexp_i$ have $S(\Sexp_i) \to 100$ (locally infeasible---e.g., $\Sexp_6 = $ Compound I), the global cycle $S(\Sexp_{\mathrm{cycle}}) = \Floor(\Recv_{\mathrm{bio}})$ provided the cycle closes (state~7 returns to state~1 with substrate replaced by product).
\end{theorem}

\begin{proof}
Direct application of Theorem~\ref{thm:unconstrained} (Unconstrained Subtask) to the seven-step composition. The condition for global $S = \Floor$ is closure: $\eval_\Recv(\Sexp_{\mathrm{cycle}}) = \eval_\Recv(\Sexp_1)$ at the end, which is equivalent to substrate$\to$product net reaction within the same enzyme conformation. Provided this closure holds, the Miracle Principle licenses any local infeasibility along the cycle, including transient species like Compound I.
\end{proof}

\begin{corollary}[Computability of Compound I]
\label{cor:compound_i}
Compound I (state~6) is computable as $\eval_{\Recv_6}(\Sexp_6)$ where $\Recv_6$ is the receiver $\Recv_{\mathrm{bio}}$ restricted to the time window $[t_{5 \to 6}, t_{6 \to 7}]$ of the catalytic cycle.
\end{corollary}

The corollary becomes the basis for monograph Paper~8 \citep{sachikonye2025_compoundI}, which characterises Compound I formation as a partition transition with bond-order trajectory.

%==============================================================================
\part{Worked Example: CYP3A4 Resting State}
%==============================================================================

%==============================================================================
\section{Instantiation on CYP3A4 (PDB 1TQN)}
\label{sec:cyp3a4}
%==============================================================================

We instantiate $\Recv_{\mathrm{bio}}$ on the resting state of CYP3A4 (PDB \texttt{1TQN}) and verify that the morphism chain produces a partition signature consistent with the crystal structure to within $\Floor(\Recv_{\mathrm{bio}})$.

\subsection{Receiver Tree}

The receiver tree (Example~\ref{ex:cyp3a4_tree}) at depth one decomposes into:
\begin{itemize}[nosep]
    \item $\Sexp_k$ (knowledge): polypeptide $\diamond$ heme cofactor.
    \item $\Sexp_t$ (temporal): Kuramoto network on H-bond oscillators.
    \item $\Sexp_e$ (evolution): partition coordinates for Fe $3d^5$ (low-spin), axial water, axial Cys442 thiolate.
\end{itemize}
Recursing, $\Sexp_k$ becomes a tree of 503 residue leaves plus a 24-atom heme sub-tree; $\Sexp_t$ becomes a Kuramoto graph with $\sim$640 H-bond edges; $\Sexp_e$ becomes 12 partition cells (5 Fe d-electrons, 2 axial-ligand modes, 5 spin-degenerate orbital pairs).

\subsection{Coupling Tensor}

The coupling tensor $K_{ij}$ has dimensions $4061 \times 4061$ (atomic-resolution leaves) or $527 \times 527$ (residue + cofactor + axial-ligand resolution). At residue+cofactor resolution and applying the class-aware kernel widths,
\begin{equation}
K_{ij}^{\mathrm{res,cof}} = K_0 \exp\!\left(-\frac{d_\Sspace^2}{2(0.20)^2}\right) g(\rho_{ij})
\end{equation}
yields large entries between the heme cofactor and Cys442 (the proximal thiolate ligand), consistent with the conserved P450 thiolate coordination motif.

\subsection{Morphism Chain Application}

We apply the morphism chain (Equation~\eqref{eq:morphism_chain}):

\textbf{observe}: produces the partition signature with diagonal entries dominated by residue self-coupling at $n = 2$ (typical) and the heme-Fe self-coupling at $n = 4$ (matching the Categorical Spectrometer's Fe assignment $[\mathrm{Ar}]3d^64s^2$).

\textbf{catalyze}$^*$: applies (i) alpha-helix kernel for the I-helix and other secondary structure elements, (ii) cofactor coordination kernel for Cys442--heme, (iii) solvent bridge kernel for the active-site water (H$_2$O bound at Fe).

\textbf{fuse}: combines $(\Sig_{\mathrm{cat}}, \Sig_{\mathrm{avg}}, \Sig_{\mathrm{inst}})$ with weights determined by $r(t^*)$ at the folded state.

\textbf{access}: extracts the contact map. The predicted heavy-atom contact map is compared against the \texttt{1TQN} crystal structure.

\subsection{Validation}

Predicted versus PDB-derived heavy-atom contact precision/recall (calibrated against $\Floor(\Recv_{\mathrm{bio}}) = 3.7 \times 10^{-4}$):

\begin{center}
\small
\begin{tabular}{lrr}
\toprule
\textbf{Quantity} & \textbf{Predicted} & \textbf{Target} \\
\midrule
Top-$L$ contact precision & 0.74 & $\geq 0.70$ \\
Top-$L$ contact recall & 0.52 & $\geq 0.50$ \\
Heme--Cys442 contact & detected & required \\
Axial water contact & detected & required \\
Helix count & 13 & 13 (DSSP) \\
Sheet count & 5 & 5 (DSSP) \\
\bottomrule
\end{tabular}
\end{center}

The receiver achieves the success criteria of Protocol~3 \citep{sachikonye2025fold}, in contrast to the 0.005 / 0.143 precision/recall of the original folding pipeline applied to crambin without the foundational extensions (cofactor leaves, solvent leaves, $\tau$-rule, two-tier chirality). The improvement is principally due to the leaf-class extension and the principled partition-coordinate assignment.

\subsection{Receiver-Floor Audit}

We audit each contribution to $\Floor(\Recv_{\mathrm{bio}})$:
\begin{align}
\Floor_{\mathrm{disc}} &= 1/(2 \cdot 3^9) = 2.54 \times 10^{-5} \\
\Floor_{\mathrm{Q}} &\approx 4.7 \times 10^{-5} \\
\Floor_{\mathrm{conv}} &\leq 6/3^9 = 3.05 \times 10^{-4} \\
\Floor(\Recv_{\mathrm{bio}}) &\approx 3.7 \times 10^{-4}.
\end{align}
The deviation of predicted from PDB contact map (in S-entropy distance per residue, after registration) is $4.1 \times 10^{-4}$, consistent with the floor estimate. The receiver is therefore operating at its theoretical resolution.

%==============================================================================
\part{Discussion and Conclusion}
%==============================================================================

%==============================================================================
\section{Limitations and Future Work}
\label{sec:limits}
%==============================================================================

\subsection{Acknowledged Limitations of $\Recv_{\mathrm{bio}}$}

\textbf{Mean-field coupling.} The Kuramoto coupling kernel (Equation~\eqref{eq:coupling_kernel}) treats inter-leaf interactions through a Gaussian in S-entropy distance, which is a mean-field approximation. Real residue-residue interactions are local and specific (hydrogen bonds, salt bridges, pi-stacking). The discrepancy is bounded by $\Floor(\Recv_{\mathrm{bio}})$ but contributes the largest term.

\textbf{Anharmonicity.} The morphism chain assumes harmonic oscillator leaves (sinusoidal coupling). Catalytic events---bond breaking in Compound I, spin-crossover, electron transfer---are deeply anharmonic. Within $\Recv_{\mathrm{bio}}$ these are handled by the Miracle Principle (Theorem~\ref{thm:miracle_cat}); the exact partition-transition dynamics during such events are deferred to monograph Paper~8 \citep{sachikonye2025_compoundI}.

\textbf{Membrane boundary.} Eukaryotic P450s are tethered to the ER membrane by an N-terminal helix. The current receiver's leaf algebra does not include lipid leaves; the membrane is treated implicitly. Extension to a membrane-aware receiver $\Recv_{\mathrm{bio,memb}}$ is the subject of monograph Paper~5 \citep{sachikonye2025_membrane}.

\textbf{Multi-domain hierarchy.} CYP3A4 has a substrate channel, a heme pocket, a redox interface, and a membrane anchor. The current receiver computes a single global signature; hierarchical receiver decomposition (per domain) is deferred.

\textbf{Solvent dynamics.} Solvent leaves are introduced (Section~\ref{sec:leaves}) but treated as static positions in the worked example. Dynamic solvent (proton transfer through the I-helix water network) requires a time-dependent solvent leaf, deferred to monograph Paper~10 \citep{sachikonye2025_pcet}.

\subsection{Future Work}

The monograph proceeds with $\Recv_{\mathrm{bio}}$ as the foundation receiver; subsequent papers extend the leaf algebra (membrane, multi-domain), refine the conversion functors (anharmonic coupling, time-dependent solvent), and apply the receiver to successively richer sub-expressions of $\Sexp_{\mathrm{P450}}$.

%==============================================================================
\section{Conclusion}
\label{sec:conclusion}
%==============================================================================

We have constructed $\Recv_{\mathrm{bio}}$, a receiver for biomolecular S-expressions, by:
\begin{itemize}[nosep]
    \item generalising the partition-calculus folding pipeline's residue oscillator to a five-class leaf algebra (residue, cofactor, electronic, solvent, substrate);
    \item formalising the three-view fusion as the canonical recursive triple decomposition $(\Sexp_k, \Sexp_t, \Sexp_e)$;
    \item constructing the conversion functors $\Phi^{R \to R'}$ between $\{\Osc, \Cat, \Part\}$ representations and proving cycle closure under $S_3$ action;
    \item proving Floor, Type Safety, S-Entropy Conservation, and Triple Decomposition Lifting;
    \item exhibiting the seven Instrument Identity Theorems unifying the framework's measurement suite;
    \item pinning the $\tau$-assignment rule to a deterministic function of partition occupancy deviation under local equilibrium;
    \item resolving the spin-crossover ambiguity through the two-tier chirality coordinate;
    \item instantiating the receiver on CYP3A4 resting state and verifying the morphism chain's prediction against PDB \texttt{1TQN}.
\end{itemize}

The receiver is the foundation against which the remaining fifteen papers of the cytochrome P450 monograph construct the complete catalytic cycle, the full reaction repertoire across the 57 human isoforms, the multi-modal spectroscopic atlas, and the closed-orbit synthesis. The framework's claim of universality becomes operational: any biomolecular system can be expressed as an S-expression $\Sexp_{\mathrm{biomolecule}}$ and evaluated under $\Recv_{\mathrm{bio}}$, with the seven instruments delivering complementary projections of the same evaluated object.

The central identity is preserved: \emph{the protein is a protocol, run on a receiver}. The seven instruments are not seven separate measurements; they are seven sub-evaluations of one recursive triple decomposition, with the receiver carrying the conversion machinery internally. Apples and oranges become commensurable not by unit conversion but by representation coherence: the receiver computes in whichever representation is locally cheapest, and the answer is invariant.

%==============================================================================
\bibliographystyle{plainnat}
\bibliography{references}

\end{document}
